\section{The Appointed Texts}
\label{sec:appointed-texts}

The Latin below is that of the 1962 typical edition, \work{Missale Romanum}
(Typis Polyglottis Vaticanis), printed pp.~697--698, nos.~3938--3949, and, for
the Preface, p.~291, no.~1080, read on the page images of the CMAA facsimile.
Every word of every text and of the Preface's proper body already stands in
the Mass's first publication, with the decree \latin{Urbis et Orbis} of
the Sacred Congregation of Rites of 12~December 1925, in the \work{Acta
Apostolicae Sedis} 17 (1925), pp.~666--668. The two printings differ in
orthography (the 1925 page writes \latin{ejus}, \latin{Jesu}), capitals and
pointing, in the headings of the chants, and in three places of extent: the
1925 page abbreviates every conclusion to \latin{Qui tecum}, which the 1962
page prints at length; it ends the Preface's proper body at \latin{Et
ideo\dots}, which the 1962 page completes with the common closing that already
stands in the Pustet Missal of 1862; and it ends the paschal Alleluia without
its last \latin{Alleluia}. The rubrics printed with the Mass are those of
the 1960 code and are described here, not reproduced.

The English of the scriptural texts is the Douay--Rheims Bible in Bishop
Challoner's revision, which translates the Clementine Vulgate in the Missal's
own Vulgate numbering. It is given at the canonical verses: where a chant takes
part of a verse, the Douay words for that part are given and an ellipsis marks
what the Missal leaves out, and each place where the Missal's form departs
from the Vulgate is named beneath the text. The English of the three orations
and the Preface is that of the Lasance \work{New Roman Missal} in the Benziger
revision of 1945, pp.~1257--1260 and 813--814, which prints this Mass in Latin
and English under the rubrics before 1960 (a commemoration of the occurring
Sunday after each oration, and the rank of a double of the first class). None
of this English is an approved liturgical translation of the 1962 Missal.
Psalms are cited in the Missal's numbering, with the Hebrew numbering in
parentheses.

\subsection{Introit}
\label{proper-introit}

\properlocus{\latin{Antiphona ad Introitum}, printed citation \latin{Apoc.\ 5, 12; 1, 6}; verse \latin{Ps.\ 71, 1} (Heb.\ 72). Missal p.~697, no.~3938.}

\begin{latinproper}
Dignus est Agnus, qui occísus est, accípere virtútem, et divinitátem, et
sapiéntiam, et fortitúdinem, et honórem. Ipsi glória et impérium in sǽcula
sæculórum. \textnormal{Ps.\ 71, 1} Deus, iudícium tuum Regi da: et iustítiam
tuam Fílio Regis. ℣.~Glória Patri.
\end{latinproper}

\begin{englishwitness}{Douay--Rheims (Challoner), Apoc 5:12 and 1:6}
\dots\ The Lamb that was slain is worthy to receive power and divinity and
wisdom and strength and honour \dots\ To him be glory and empire for ever and
ever. \dots
\end{englishwitness}

\begin{englishwitness}{Douay--Rheims (Challoner), Ps 71:2, the verse}
Give to the king thy judgment, O God, and to the king's son thy justice
\dots
\end{englishwitness}

\noindent The antiphon joins two verses of the Apocalypse. From Apoc 5:12 it
takes the angels' acclamation without its opening, ``Saying with a loud
voice,'' and without its last two gifts, ``and glory and benediction''; from
Apoc 1:6 it takes only the doxology, without ``And hath made us a kingdom, and
priests to God and his Father'' and without the ``Amen.'' The Missal's
\latin{Ps.\ 71, 1} counts the psalm's title, \latin{Psalmus, in Salomonem} (``A
psalm on Solomon''), as unnumbered; the Clementine Vulgate and the Douay number
the title as verse~1 and print the Introit's words at the head of their verse~2,
which goes on, ``To judge thy people with justice, and thy poor with
judgment.'' After the \latin{Gloria Patri} the antiphon is repeated (RGMR 427).

\subsection{Collect}
\label{proper-collect}

\properlocus{\latin{Oratio}. Missal p.~697, no.~3939.}

\begin{latinproper}
Omnípotens sempitérne Deus, qui in dilécto Fílio tuo, universórum Rege, ómnia
instauráre voluísti: concéde propítius; ut cunctæ famíliæ géntium, peccáti
vúlnere disgregátæ, eius suavíssimo subdántur império: Qui tecum vivit et
regnat in unitáte.
\end{latinproper}

\begin{englishwitness}{Lasance, The New Roman Missal (1945), Feast of Our Lord Jesus Christ, King, Collect}
Almighty and eternal God, Who hast wished to restore all things through Thy
beloved Son, the King of the universe, graciously grant that all the familes of
the Gentiles separated by the wound of sin, may be subjected to His most loving
dominion, Who with Thee liveth.
\end{englishwitness}

\noindent The Lasance page prints ``familes'' for ``families.'' The prayer is
addressed to the Father and names the Son at its end, so it concludes
\latin{Qui tecum} (RG 115~c). The Missal then directs that no commemoration of
the Sunday is made, and repeats the direction after the Secret and the
Postcommunion.

\subsection{Epistle}
\label{proper-epistle}

\properlocus{\latin{Lectio Epistolae beati Pauli Apostoli ad Colossenses}, Col 1:12--20. Missal pp.~697--698, no.~3940.}

\begin{latinproper}
Fratres: Grátias ágimus Deo Patri, qui dignos nos fecit in partem sortis
sanctórum in lúmine, qui erípuit nos de potestáte tenebrárum, et tránstulit
in regnum Fílii dilectiónis suæ, in quo habémus redemptiónem per sánguinem
eius remissiónem peccatórum: qui est imágo Dei invisíbilis, primogénitus
omnis creatúræ: quóniam in ipso cóndita sunt univérsa in cælis et in terra,
visibília et invisibília, sive throni, sive dominatiónes, sive principátus,
sive potestátes: ómnia per ipsum et in ipso creáta sunt: et ipse est ante
omnes, et ómnia in ipso constant. Et ipse est caput córporis Ecclésiæ, qui
est princípium, primogénitus ex mórtuis: ut sit in ómnibus ipse primátum
tenens: quia in ipso complácuit omnem plenitúdinem inhabitáre: et per eum
reconciliáre ómnia in ipsum, pacíficans per sánguinem crucis eius, sive quæ
in terris, sive quæ in cælis sunt, in Christo Iesu Dómino nostro.
\end{latinproper}

\begin{englishwitness}{Douay--Rheims (Challoner), Col 1:12--20}
Giving thanks to God the Father, who hath made us worthy to be partakers of
the lot of the saints in light: Who hath delivered us from the power of
darkness and hath translated us into the kingdom of the Son of his love, In
whom we have redemption through his blood, the remission of sins: Who is the
image of the invisible God, the firstborn of every creature: For in him were
all things created in heaven and on earth, visible and invisible, whether
thrones, or dominations, or principalities, or powers. All things were created
by him and in him. And he is before all: and by him all things consist. And he
is the head of the body, the church: who is the beginning, the firstborn from
the dead, that in all things he may hold the primacy: Because in him, it hath
well pleased the Father that all fulness should dwell: And through him to
reconcile all things unto himself, making peace through the blood of his
cross, both as to the things that are on earth and the things that are in
heaven.
\end{englishwitness}

\noindent The Missal begins with the liturgical address \latin{Fratres:
Gratias agimus}, where the Vulgate's verse 12 has the participle \latin{gratias
agentes} that the Douay renders ``Giving thanks,'' and it closes with \latin{in
Christo Iesu Domino nostro}, which the Vulgate's verse 20 does not have; the
Lasance prints these frames as ``Brethren, we give thanks to God the Father''
and ``in Christ Jesus Our Lord.'' Every other word of the nine verses agrees
with the Clementine.

\subsection{Gradual}
\label{proper-gradual}

\properlocus{\latin{Graduale}, \latin{Ps.\ 71, 8 et 11} (Heb.\ 72). Missal p.~698, no.~3941.}

\begin{latinproper}
Dominábitur a mari usque ad mare: et a flúmine usque ad términos orbis
terrárum. ℣.~Et adorábunt eum omnes reges terræ: omnes gentes sérvient ei.
\end{latinproper}

\begin{englishwitness}{Douay--Rheims (Challoner), Ps 71:8 and 11}
And he shall rule from sea to sea, and from the river unto the ends of the
earth. \dots\ And all kings of the earth shall adore him: all nations shall
serve him.
\end{englishwitness}

\noindent The respond begins \latin{Dominabitur} without the Vulgate's initial
\latin{Et}, which the Douay renders ``And''; the verse is verse~11 whole.

\subsection{Alleluia}
\label{proper-alleluia}

\properlocus{\latin{Alleluia, alleluia}, verse \latin{Dan.\ 7, 14}. Missal p.~698, no.~3942.}

\begin{latinproper}
Allelúia, allelúia. ℣.~Potéstas eius, potéstas ætérna, quæ non auferétur: et
regnum eius, quod non corrumpétur. Allelúia.
\end{latinproper}

\begin{englishwitness}{Douay--Rheims (Challoner), Dan 7:14}
\dots\ his power is an everlasting power that shall not be taken away: and his
kingdom that shall not be destroyed.
\end{englishwitness}

\noindent The verse is the second half of Dan 7:14 only. The first half,
``And he gave him power, and glory, and a kingdom: and all peoples, tribes, and
tongues shall serve him,'' is not sung. The Mass has no Sequence.

\subsection{Gospel}
\label{proper-gospel}

\properlocus{\latin{Sequentia sancti Evangelii secundum Ioannem}, Jn 18:33--37. Missal p.~698, no.~3945.}

\begin{latinproper}
In illo témpore: Dixit Pilátus ad Iesum: Tu es Rex Iudæórum? Respóndit Iesus:
A temetípso hoc dicis, an álii dixérunt tibi de me? Respóndit Pilátus: Numquid
ego Iudǽus sum? Gens tua et pontífices tradidérunt te mihi: quid fecísti?
Respóndit Iesus: Regnum meum non est de hoc mundo. Si ex hoc mundo esset regnum
meum, minístri mei útique decertárent, ut non tráderer Iudǽis: nunc autem regnum
meum non est hinc. Dixit ítaque ei Pilátus: Ergo Rex es tu? Respóndit Iesus: Tu
dicis, quia Rex sum ego. Ego in hoc natus sum, et ad hoc veni in mundum, ut
testimónium perhíbeam veritáti: omnis qui est ex veritáte, audit vocem meam.
\end{latinproper}

\begin{englishwitness}{Douay--Rheims (Challoner), Jn 18:33--37}
\dots\ Art thou the king of the Jews? Jesus answered: Sayest thou this thing
of thyself, or have others told it thee of me? Pilate answered: Am I a Jew?
Thy own nation and the chief priests have delivered thee up to me. What hast
thou done? Jesus answered: My kingdom is not of this world. If my kingdom were
of this world, my servants would certainly strive that I should not be
delivered to the Jews: but now my kingdom is not from hence. Pilate therefore
said to him: Art thou a king then? Jesus answered: Thou sayest that I am a
king. For this was I born, and for this came I into the world; that I should
give testimony to the truth. Every one that is of the truth heareth my voice.
\end{englishwitness}

\noindent The liturgical opening \latin{In illo tempore: Dixit Pilatus ad
Iesum}, which the Lasance renders ``At that time: Pilate said to Jesus,''
stands for the Vulgate's \latin{Introivit ergo iterum in praetorium Pilatus:
et vocavit Jesum, et dixit ei}, ``Pilate therefore went into the hall again
and called Jesus and said to him.'' Every other word of the five verses agrees
with the Clementine. The Credo follows.

\subsection{Offertory}
\label{proper-offertory}

\properlocus{\latin{Antiphona ad Offertorium}, \latin{Ps.\ 2, 8}. Missal p.~698, no.~3946.}

\begin{latinproper}
Póstula a me, et dabo tibi gentes hereditátem tuam, et possessiónem tuam
términos terræ.
\end{latinproper}

\begin{englishwitness}{Douay--Rheims (Challoner), Ps 2:8}
Ask of me, and I will give thee the Gentiles for thy inheritance, and the
utmost parts of the earth for thy possession.
\end{englishwitness}

\subsection{Secret}
\label{proper-secret}

\properlocus{\latin{Secreta}. Missal p.~698, no.~3947.}

\begin{latinproper}
Hóstiam tibi, Dómine, humánæ reconciliatiónis offérimus: præsta, quǽsumus;
ut, quem sacrifíciis præséntibus immolámus, ipse cunctis géntibus unitátis et
pacis dona concédat, Iesus Christus, Fílius tuus, Dóminus noster: Qui tecum
vivit et regnat in unitáte.
\end{latinproper}

\begin{englishwitness}{Lasance, The New Roman Missal (1945), Feast of Our Lord Jesus Christ, King, Secret}
O Lord, we offer Thee this host for the reconciliation of humanity; grant, we
beseech Thee, that Jesus Christ Thy Son our Lord, Whom we immolate in this
sacrifice, will bestow on all Gentiles the gifts of unity and peace, Who with
Thee liveth.
\end{englishwitness}

\noindent The Secret is said secretly; the Missal then directs the proper
Preface.

\subsection{Preface}
\label{proper-preface}

\properlocus{\latin{Praefatio de D.~N. Iesu Christo Rege}, among the \latin{Praefationes sine cantu}. Missal p.~291, no.~1080; RGMR 492.}

\begin{latinproper}
Vere dignum et iustum est, æquum et salutáre, nos tibi semper et ubíque
grátias ágere: Dómine, sancte Pater, omnípotens ætérne Deus: Qui unigénitum
Fílium tuum, Dóminum nostrum Iesum Christum, Sacerdótem ætérnum et
universórum Regem, óleo exsultatiónis unxísti: ut, seípsum in ara Crucis
hóstiam immaculátam et pacíficam ófferens, redemptiónis humánæ sacraménta
perágeret: et, suo subiéctis império ómnibus creatúris, ætérnum et universále
regnum imménsæ tuæ tráderet maiestáti: regnum veritátis et vitæ; regnum
sanctitátis et grátiæ; regnum iustítiæ, amóris et pacis. Et ídeo cum Angelis
et Archángelis, cum Thronis et Dominatiónibus, cumque omni milítia cæléstis
exércitus, hymnum glóriæ tuæ cánimus, sine fine dicéntes:
\end{latinproper}

\begin{englishwitness}{Lasance, The New Roman Missal (1945), Preface of Jesus Christ, King}
It is truly meet and just, right and profitable, for us, at all times, and in
all places, to give thanks to Thee, O Lord, the holy One, the Father almighty,
the everlasting God: Who didst anoint with the oil of gladness Thine
only-begotten Son, our Lord Jesus Christ, eternal Priest and King of the
universe: that, offering Himself a stainless peace-offering on the altar of
the cross, He might fulfill the pledges of man's redemption; and, having all
creatures subject to His power, might deliver to Thy sublime majesty an
eternal and universal kingdom, a kingdom of truth and life; a kingdom of
holiness and grace, a kingdom of justice, love and peace. And therefore with
the angels and archangels, with the thrones and dominations, and with all the
array of the heavenly host, we sing a hymn to Thy glory and unceasingly
repeat:
\end{englishwitness}

\noindent The \latin{Sanctus} follows, and the dialogue before the Preface is
that of the Ordinary. The Missal prints this Preface among the Prefaces, with
the rubric that it is said in festive and votive Masses of Christ the King;
the 1925 page had printed it in place, between the Secret and the Communion.
The two printings group the opening vocatives differently, \latin{Domine
sancte, Pater omnipotens} in 1925 and \latin{Domine, sancte Pater,
omnipotens} in 1962, without changing a word.

\subsection{Communion}
\label{proper-communion}

\properlocus{\latin{Antiphona ad Communionem}, \latin{Ps.\ 28, 10 et 11} (Heb.\ 29). Missal p.~698, no.~3948.}

\begin{latinproper}
Sedébit Dóminus Rex in ætérnum: Dóminus benedícet pópulo suo in pace.
\end{latinproper}

\begin{englishwitness}{Douay--Rheims (Challoner), Ps 28:10--11}
\dots\ and the Lord shall sit king for ever. \dots\ the Lord will bless his
people with peace.
\end{englishwitness}

\noindent The antiphon joins the second half of verse~10 to the second half of
verse~11. It leaves out \latin{Dominus diluvium inhabitare facit}, ``The Lord
maketh the flood to dwell,'' and \latin{Dominus virtutem populo suo dabit},
``The Lord will give strength to his people.'' The Challoner text prints the
Vulgate's verses 10 and 11 as one verse numbered~10, so its Psalm~28 has ten
verses where the Missal and the Clementine count eleven.

\subsection{Postcommunion}
\label{proper-postcommunion}

\properlocus{\latin{Postcommunio}. Missal p.~698, no.~3949.}

\begin{latinproper}
Immortalitátis alimóniam consecúti, quǽsumus, Dómine: ut, qui sub Christi
Regis vexíllis militáre gloriámur, cum ipso, in cælésti sede, iúgiter regnáre
possímus: Qui tecum vivit et regnat in unitáte.
\end{latinproper}

\begin{englishwitness}{Lasance, The New Roman Missal (1945), Feast of Our Lord Jesus Christ, King, Postcommunion}
Fed with this immortal nourishment, we beseech Thee, O Lord, that we who glory
to fight under the standard of Christ the King, may forever reign with Him on
the heavenly throne. Who with Thee liveth.
\end{englishwitness}

\noindent The Latin has the plural \latin{vexillis}, the standards; the
Lasance renders it in the singular. The Missal prints for this Mass no Sequence, proper
\latin{Communicantes} or \latin{Hanc igitur}, prayer over the people, second
oration or alternative reading.

\subsection{Printed for Other Days: The Tract and the Paschal Alleluia}
\label{sec:other-days}

Between the Alleluia and the Gospel the Missal prints two chants that replace
the Gradual and Alleluia in other seasons. The first stands under the rubric
\latin{In Missis votivis post Septuagesimam, omissis Alleluia et versu
sequenti, dicitur}: in votive Masses after Septuagesima, the Alleluia and its
verse are left out and this Tract is said (no.~3943, \latin{Ps.\ 88, 27-28 et
30}). The second stands under \latin{Tempore autem paschali omittitur
graduale, et eius loco dicitur}: in Paschaltide the Gradual is left out and
this Alleluia is said in its place (no.~3944). The feast is kept on the last
Sunday of October, which lies neither between Septuagesima and Easter nor in
Paschaltide, so neither chant is ever sung on the feast. Both belong to the
votive Mass of Christ the King, the same Mass said on other days (RGMR
308~a~5).

\begin{latinproper}
\textnormal{Tractus.} Ipse invocábit me: Pater meus es tu: Deus meus, et
suscéptor salútis meæ. ℣.~Et ego primogénitum ponam illum, excélsum præ
régibus terræ. ℣.~Et ponam in sǽculum sǽculi semen eius: et thronum eius
sicut dies cæli.
\end{latinproper}

\begin{englishwitness}{Douay--Rheims (Challoner), Ps 88:27--28 and 30}
He shall cry out to me: Thou art my father: my God, and the support of my
salvation. And I will make him my firstborn, high above the kings of the
earth. \dots\ And I will make his seed to endure for evermore: and his throne
as the days of heaven.
\end{englishwitness}

\begin{latinproper}
\textnormal{Alleluia in Paschaltide.} Allelúia, allelúia. ℣.~Potéstas eius,
potéstas ætérna, quæ non auferétur: et regnum eius, quod non corrumpétur.
Allelúia. ℣.~\textnormal{Apoc.\ 19, 16} Habet in vestiménto et in fémore suo
scriptum: Rex regum et Dóminus dominántium. Allelúia.
\end{latinproper}

\begin{englishwitness}{Douay--Rheims (Challoner), Apoc 19:16, the second verse}
And he hath on his garment and on his thigh written: KING OF KINGS AND LORD OF
LORDS.
\end{englishwitness}

\noindent The Tract takes verses 27, 28 and 30 of Psalm~88 (89), leaving out
verse~29; the paschal Alleluia's first verse is the Alleluia verse of the
feast, and its second is Apoc 19:16 without the Vulgate's initial \latin{Et},
which the Douay renders ``And.''
